\documentclass[10pt]{article}

\usepackage{fancyhdr}
\pagestyle{fancy}
\fancyhf{}
\fancyhead[C]{Working Note: Consolidation, Proofs Checked Numerically Only}
\fancyfoot[C]{\thepage}
\renewcommand{\headrulewidth}{0.4pt}

\usepackage[margin=2cm]{geometry}
\usepackage[T1]{fontenc}
\usepackage[utf8]{inputenc}
\usepackage{amsmath,amssymb,amsthm,mathtools}
\usepackage{booktabs}
\usepackage{tabularx}
\usepackage{enumitem}
\usepackage{xcolor}
\usepackage[colorlinks=true,linkcolor=blue!55!black,urlcolor=blue!55!black]{hyperref}

\setlist{nosep,leftmargin=1.3em}

\theoremstyle{plain}
\newtheorem{proposition}{Proposition}
\newtheorem{corollary}{Corollary}
\theoremstyle{remark}
\newtheorem{remark}{Remark}

\newcommand{\occ}{d^{\pi}}
\newcommand{\E}{\mathbb E}
\newcommand{\Scal}{\mathcal S}
\newcommand{\Acal}{\mathcal A}
\newcommand{\Zcal}{\mathcal Z}
\newcommand{\Ccal}{\mathcal C}
\newcommand{\dd}{\mathrm d}
\newcommand{\open}[1]{\textcolor{red!70!black}{[Open: #1]}}

\title{The error-controlled regulator is a cause-controlled regulator on its own model of the cause}
\author{Working note for the RLC draft, 2026-11-04}
\date{}

\begin{document}
\maketitle

\section{The frame}
An error-controlled regulator reads effects, infers a cause, and acts on the inferred cause as if it had read it.
It therefore splits into two maps, inference from histories to beliefs and control from beliefs to actions, and the good regulator theorem applies to their composition.
The belief is the regulator's own event set, the set $R$ of Conant and Ashby, so the mapping $h\colon S \to R$ of their theorem sends system states to belief states and the belief is the literal object of the aphorism.
The policy acts on the whole belief and not on a point estimate of the cause, so the right object is cause control on the belief, never on a predicted value.

A disturbance is whatever enters between what the regulator reads and what it controls, and the same structure recurs one level up.
Each level has two floors of different kinds, what it cannot read and what it reads but cannot counter.

\begin{center}
\footnotesize
\begin{tabularx}{\linewidth}{lXXXXX}
\toprule
level & reads & infers & acts on & cannot read & cannot counter \\
\midrule
C\&A cause control & $D$ & nothing & $D$ & nothing & variety of $A$ against $D$ \\
policy in the POMDP & $X$ & belief over $S$ and $\theta$ & belief & $\varepsilon$, the noise term & variety of $A$ against the inferred cause \\
learner & outcomes, own updates & law of the policy, drift & policy parameters & exogenous drift, through effects only & capacity of the policy class \\
\bottomrule
\end{tabularx}
\end{center}

The level above is cause-controlled on what it can read before it takes effect, which is why endogenous disturbances admit cause control and exogenous ones do not.
Requisite variety bounds the model from below, since the representation must keep the distinctions in the inferred cause that the actions have to counter.
Minimality bounds it from above, since nothing of the cause is kept beyond what changes the action.
The simplest optimal regulator sits on the lower bound, which turns Conant and Ashby's simplest into a representation criterion.

The two statements below make the policy row exact.
The first splits the noise term of Theorem~1 of the draft into what a cause reader would still face and the price of inferring the cause instead of reading it.
The second is requisite variety with the noise floor and the invariants of $\psi^{-1}$ accounted for.

\section{Setting}
The cause is $C = (S, \theta)$, the hidden state and the law, with the regular part of the disturbance folded into it.
The regulator reads the history $X$, holds the belief $b = b(X)$ over $C$, and chooses $A$ from $X$ and independent randomness.
The irregular part $\varepsilon$ of the disturbance is what no history reveals, and the outcome is $Z = \psi((S,\varepsilon), A)$ with channel $\kappa(c,a)$ the law of $Z$ given $(C,A) = (c,a)$.
The one structural assumption is that $\varepsilon$ is drawn independently of the past given the pair, so that
\begin{equation}
  Z \perp X \mid (C, A).
  \label{eq:irregular}
\end{equation}
The observable class is $\hat Z = \E[\kappa(C,A) \mid X, A]$, which by \eqref{eq:irregular} equals $\bar\kappa(b, A)$ with $\bar\kappa(b,a) = \int \kappa(c,a)\, b(\dd c)$, the channel of the belief MDP.
Bases are $\lambda$ on $\Zcal$, a reference on $\Ccal$ where needed, and counting on finite sets, and $\ell = \ln \tfrac{\dd\nu}{\dd\lambda}$.

\section{The price of inferring the cause}
\begin{proposition}[Noise split]
\label{prop:noise-split}
For any history policy, under \eqref{eq:irregular} and finite entropies,
\begin{equation}
  H_\nu(Z) = \underbrace{H_\lambda(Z \mid C, A)}_{\text{irreducible noise}} + \underbrace{I(C; Z \mid X, A)}_{\text{hidden cause}} + \underbrace{I(\hat Z; Z)}_{\text{control net of ambiguity}} + \E[\ell(Z)],
  \label{eq:noise-split}
\end{equation}
and the noise term of Theorem~1 is $H_\lambda(Z \mid \hat Z) = H_\lambda(Z \mid X, A) = H_\lambda(Z \mid b, A) = H_\lambda(Z \mid C,A) + I(C; Z \mid X, A)$.
\end{proposition}
\begin{proof}
The chain rule gives $H_\lambda(Z) = H_\lambda(Z \mid X,A) + I(X,A; Z)$.
Since $\hat Z$ is the law of $Z$ given $(X,A)$, the pair $(X,A)$ is independent of $Z$ given $\hat Z$, so $H_\lambda(Z \mid X,A) = H_\lambda(Z \mid \hat Z)$ and $I(X,A;Z) = I(\hat Z; Z)$.
Since $\hat Z$ is a function of $(b,A)$ and $(b,A)$ is a function of $(X,A)$, also $H_\lambda(Z \mid b,A) = H_\lambda(Z \mid X,A)$.
The chain rule on $C$ gives $H_\lambda(Z \mid X,A) = H_\lambda(Z \mid X,A,C) + I(C;Z \mid X,A)$, and \eqref{eq:irregular} gives $H_\lambda(Z \mid X,A,C) = H_\lambda(Z \mid C,A)$.
Part~1 of the alignment proposition adds $\E[\ell(Z)]$ to pass from $\lambda$ to $\nu$.
\end{proof}

\begin{remark}
A cause-controlled regulator reads $C$ before acting, so for it $X$ contains $C$ and the hidden-cause term is zero.
The error-controlled regulator pays $I(C; Z \mid X, A)$ in the noise term for inferring the cause instead of reading it, and the belief MDP treats this price as channel noise of $\bar\kappa$.
This is the exact form of the sentence in paragraph~3 of the draft, that in general the noise term also holds the hidden state and the unknown law.
\end{remark}

\begin{remark}
With $\tilde Z = \kappa(C,A)$ the true class, $Z$ depends on $C$ only through $\tilde Z$ given $A$, so by data processing $I(C; Z \mid X, A) \le I(C; \tilde Z \mid X, A) = \bar I_P + \bar I_\kappa + \bar I_S$ of the split proposition, with gap $I(C; \tilde Z \mid X, A, Z)$.
The hidden-cause term of \eqref{eq:noise-split} is therefore the part of those three floors that the outcome draw itself reveals.
\end{remark}

\section{Requisite variety with noise and invariants}
\begin{proposition}[Requisite variety]
\label{prop:variety}
For finite $\Ccal$, $\Acal$ and $\Zcal$ with counting bases,
\begin{equation}
  H(Z) \;\ge\; \big[H(C) - H(C \mid Z, A)\big] - H(A) + H(Z \mid C, A).
  \label{eq:variety}
\end{equation}
For a deterministic policy on a finite belief alphabet, $A = \pi(b)$ gives $H(A) \le H(b)$, hence
\begin{equation}
  H(b) \;\ge\; \big[H(C) - H(C \mid Z, A)\big] - H(Z) + H(Z \mid C, A).
  \label{eq:variety-model}
\end{equation}
\end{proposition}
\begin{proof}
$H(Z) \ge H(Z \mid A) = H(C, Z \mid A) - H(C \mid Z, A) = H(C \mid A) + H(Z \mid C, A) - H(C \mid Z, A)$, and $H(C \mid A) = H(C) - I(C;A) \ge H(C) - H(A)$.
For $A = \pi(b)$ the data processing inequality gives $H(A) \le H(b)$.
\end{proof}

\begin{remark}
The bracket $H(C) - H(C \mid Z, A) = I(C; Z, A)$ is the variety of the cause that the outcome and the action together reflect.
The invariant part $H(C \mid Z, A)$, the distinctions in the cause that $\psi$ does not pass to the outcome under any action, is exempt from the bound, which is the pre-image $\psi^{-1}(z)$ of the walk read as a discount on requisite variety.
Ashby's form $H(Z) \ge H(D) - H(R)$ is the case of no invariants and no noise.
\end{remark}

\begin{remark}
Inequality \eqref{eq:variety-model} is requisite variety on the model.
The belief must carry at least the variety of the cause that reaches the outcome, less the entropy the regulator tolerates in the outcome, plus the noise it cannot remove.
This bounds the representation from below, and minimality bounds it from above.
\end{remark}

\section{The nested regulator}
\begin{corollary}
For every history policy the outcome law equals that of the Markov policy on the belief with the same marginal laws of $(b_t, A_t)$, and by Theorem~1 the minimum over history policies is attained by a deterministic policy on the belief.
That policy is a cause-controlled regulator on the belief MDP with channel $\bar\kappa$, whose criterion value exceeds the irreducible noise of a cause reader by the hidden-cause term of \eqref{eq:noise-split}.
The class of regulators that read $C$ contains the history policies, so $\min_{\Pi} H_\nu(Z) \ge \min_{\Pi_C} H_\nu(Z)$.
\end{corollary}
The first sentence is the belief part of the proof of Theorem~1, the second is Proposition~\ref{prop:noise-split}, and the third is set inclusion.

\section{Model criteria as properties of the inference map}
Each criterion of the draft is a property of inferring the cause well enough to be cause-controlled on it.
\begin{itemize}
  \item Sufficiency, the belief carries everything the history says about the cause, which is $H_\lambda(Z \mid b,A) = H_\lambda(Z \mid X,A)$ in Proposition~\ref{prop:noise-split}.
  \item Closure, the inferred cause updates from itself and the next effect, which is the filter.
  \item Consistency, the predicted law equals the true one, so acting on the inferred cause is acting on the cause.
  \item Identifiability, the effects reveal the cause distinctions that matter, which is the bracket of \eqref{eq:variety} being attainable from $X$.
  \item Minimality, nothing of the cause is kept beyond what changes the action.
\end{itemize}
The two coarse-grainings of the supplementary are the two ends of this map, the history coarse-grained to the belief and the belief coarse-grained to the outcome laws it selects.
The goal measure $\nu$ decides which distinctions in the inferred cause matter, and the information floor says the regulator must keep at least the cause distinctions that reach $G$ while laziness drops the rest.

\section{Open points}
\begin{itemize}
  \item \open{the vertex argument of Theorem~1 needs a finite belief alphabet, and \eqref{eq:variety-model} uses counting entropy of $b$.}
  \item \open{for continuous $C$ the identities hold with product bases as in the primer, and \eqref{eq:variety} needs the base on $\Ccal$ made explicit.}
  \item \open{the learner row of the table is a frame, not a result. The claim that an endogenous disturbance admits cause control one level up needs its own statement with the learner's objective.}
  \item \open{Conant and Ashby's residual variation under error control is a temporal argument, the error must occur before it is corrected. Proposition~\ref{prop:noise-split} formalises a different mechanism, that the regulator never reads $\varepsilon$, and the draft should say so where it claims to formalise their frame.}
\end{itemize}
Both identities were checked numerically on random finite laws, script \texttt{nested\_check.py} in the session scratchpad.

\end{document}
